\section*{Capítulo \ref{ch:g5:series-parallel-circuits} --- Circuitos en serie y en paralelo}

\begin{solution}{exo:g5:series-parallel-circuits:1}
Una \omterm{def:g5:series-parallel-circuits:parallel}{rama} es una de las carreteras separadas en las que se
bifurca el circuito; un \omterm{def:g5:series-parallel-circuits:parallel}{nudo} es un punto donde las carreteras se
bifurcan o se vuelven a juntar. Dos \omterm{def:g3:first-electric-circuit:bulb}{bombillas} \omterm{def:g5:series-parallel-circuits:parallel}{en paralelo} viajan
por dos \omterm{def:g5:series-parallel-circuits:parallel}{ramas} con \omterm{def:g3:first-electric-circuit:bulb}{bombilla}.
\end{solution}

\begin{solution}{exo:g5:series-parallel-circuits:2}
Independencia (un hueco para solo su propia \omterm{def:g5:series-parallel-circuits:parallel}{rama}); brillo completo
(cada \omterm{def:g3:first-electric-circuit:bulb}{bombilla} brilla como si estuviera sola, a cambio de que la
\omterm{def:g3:first-electric-circuit:battery}{pila} se gaste más deprisa); puestos de mando (un \omterm{ex:g3:first-electric-circuit:switch}{interruptor} de
\omterm{def:g5:series-parallel-circuits:parallel}{rama} manda sobre su \omterm{def:g5:series-parallel-circuits:parallel}{rama} y un \omterm{ex:g3:first-electric-circuit:switch}{interruptor} anterior a la
bifurcación manda sobre todas).
\end{solution}

\begin{solution}{exo:g5:series-parallel-circuits:3}
La \omterm{def:g3:first-electric-circuit:bulb}{bombilla} 2 sigue encendida, sin enterarse --- su propia carretera
está completa. \omterm{def:g4:series-circuits:series}{En serie} se apagaba, porque allí las dos
\omterm{def:g3:first-electric-circuit:bulb}{bombillas} compartían una única carretera y cualquier hueco la
rompía para las dos.
\end{solution}

\begin{solution}{exo:g5:series-parallel-circuits:4}
La pareja B (\omterm{def:g5:series-parallel-circuits:parallel}{en paralelo}) brilla más por \omterm{def:g3:first-electric-circuit:bulb}{bombilla} --- cada una a
pleno brillo. La pareja B gasta también antes su \omterm{def:g3:first-electric-circuit:battery}{pila}: alimenta dos
\omterm{def:g5:series-parallel-circuits:parallel}{ramas} completas a la vez.
\end{solution}

\begin{solution}{exo:g5:series-parallel-circuits:5}
Los dos cerrados: las dos \omterm{def:g3:first-electric-circuit:bulb}{bombillas} brillan. El general abierto: las
dos apagadas --- manda sobre la carretera común. El general cerrado y
el \omterm{ex:g3:first-electric-circuit:switch}{interruptor} de \omterm{def:g5:series-parallel-circuits:parallel}{rama} abierto: la \omterm{def:g3:first-electric-circuit:bulb}{bombilla} 1 brilla y la
\omterm{def:g3:first-electric-circuit:bulb}{bombilla} 2 está apagada.
\end{solution}

\begin{solution}{exo:g5:series-parallel-circuits:6}
\omterm{def:g4:series-circuits:series}{En serie}, una \omterm{def:g3:first-electric-circuit:bulb}{bombilla} fundida (¡o una lámpara apagada!) dejaría
a oscuras la casa entera, y todo se iría atenuando según se sumaran
luces. Comodidades del montaje \omterm{def:g5:series-parallel-circuits:parallel}{en paralelo}: el \omterm{ex:g3:first-electric-circuit:switch}{interruptor} de
cada habitación manda solo sobre su habitación; una \omterm{def:g3:first-electric-circuit:bulb}{bombilla} fundida
en la cocina deja el resto de la casa iluminado; todas las lámparas
lucen a pleno brillo.
\end{solution}

\begin{solution}{exo:g5:series-parallel-circuits:7}
Antes de la bifurcación --- en la carretera común por la que la
corriente entra en la casa: el \omterm{ex:g3:first-electric-circuit:switch}{interruptor} general del cuadro de
fusibles. Corta la casa entera de golpe y se usa para reparaciones y
emergencias.
\end{solution}

\begin{solution}{exo:g5:series-parallel-circuits:8}
\omterm{def:g5:series-parallel-circuits:parallel}{En paralelo}: apagar la \omterm{def:g1:light-and-shadows:source}{luz} del techo no mata la lámpara de
mesilla, y desenchufar la lámpara no apaga el techo --- cada una viaja
por su propia \omterm{def:g5:series-parallel-circuits:parallel}{rama}.
\end{solution}

\begin{solution}{exo:g5:series-parallel-circuits:9}
Con nuestros \omterm{ex:g3:first-electric-circuit:switch}{interruptores} corrientes los diseños fracasan una y otra
vez: un \omterm{ex:g3:first-electric-circuit:switch}{interruptor} \omterm{def:g4:series-circuits:series}{en serie} con la \omterm{def:g3:first-electric-circuit:bulb}{bombilla} solo se puede
abrir desde un extremo; dos \omterm{ex:g3:first-electric-circuit:switch}{interruptores} corrientes \omterm{def:g4:series-circuits:series}{en serie}
exigen que \emph{los dos} estén cerrados (cualquier extremo puede
apagar la \omterm{def:g1:light-and-shadows:source}{luz}, pero no siempre encenderla); dos \omterm{def:g5:series-parallel-circuits:parallel}{en paralelo}
exigen que los dos estén abiertos para que haya oscuridad. La
comodidad del pasillo necesita un \omterm{ex:g3:first-electric-circuit:switch}{interruptor} que \emph{elija entre
dos carreteras} en vez de limitarse a cortar una --- el conmutador,
que los pasillos de verdad usan por parejas.
\end{solution}

\begin{solution}{exo:g5:series-parallel-circuits:10}
La comodidad: las \omterm{def:g3:first-electric-circuit:bulb}{bombillas} viajan por \omterm{def:g5:series-parallel-circuits:parallel}{ramas} paralelas, así que
una fundida solo se apaga a sí misma. El aviso: cada \omterm{def:g3:first-electric-circuit:bulb}{bombilla} que
falta quita una \omterm{def:g5:series-parallel-circuits:parallel}{rama}, y la \omterm{def:g3:first-electric-circuit:battery}{pila} (o la pequeña fuente de
alimentación de la guirnalda) empuja todo su esfuerzo por las
\omterm{def:g5:series-parallel-circuits:parallel}{ramas} que quedan --- si faltan demasiadas, cada superviviente
recibe más de lo que se construyó para aguantar, se calienta más de la
cuenta y se funde joven.
\end{solution}

\begin{solution}{exo:g5:series-parallel-circuits:11}
En el propio tramo de carretera del \omterm{ex:g3:first-electric-circuit:switch}{interruptor} de \omterm{def:g5:series-parallel-circuits:parallel}{rama}: entre ese
\omterm{ex:g3:first-electric-circuit:switch}{interruptor} y la \omterm{def:g3:first-electric-circuit:bulb}{bombilla} 2 (antes de que su \omterm{def:g5:series-parallel-circuits:parallel}{rama} se vuelva
a juntar con el resto). Ahí se enciende exactamente cuando esa
\omterm{def:g5:series-parallel-circuits:parallel}{rama} funciona --- y queda \omterm{def:g4:series-circuits:series}{en serie} con la \omterm{def:g3:first-electric-circuit:bulb}{bombilla} 2, porque
las dos viajan por la misma \omterm{def:g5:series-parallel-circuits:parallel}{rama}.
\end{solution}
